% ============================================================
% CAPITOLO 10 - IL SECONDO RIVOLTO: QUANDO E COME USARLO
% ============================================================
\chapter{Il secondo rivolto: quando e come usarlo}
\label{chap:quarta_sesta}

Il Capitolo \ref{chap:rivolti} ha mostrato come si costruisce il secondo rivolto della triade. Qui si vede perché il suo uso è limitato e in quali tre casi è invece ammesso, e anzi indispensabile.

\section{Perché il 6/4 è problematico}
\label{sec:problema_secondo_rivolto}

Il secondo rivolto è formato dalle stesse note di una triade consonante, ma viene percepito come \textbf{dissonante}. Il motivo è la quarta tra il basso e la fondamentale: l'orecchio sente il basso come fondamentale dell'accordo, mentre la struttura delle note ne indica un'altra. Questa ambiguità produce instabilità, e per questo l'accordo di quarta e sesta non si usa liberamente come gli altri rivolti (la spiegazione acustica è nell'Approfondimento \ref{phys:conflitto_6_4}).

Dall'ambiguità discende la regola pratica: il secondo rivolto va sempre preparato e risolto, e compare solo in tre contesti, trattati nelle sezioni seguenti. La sintesi è nel Capitolo \ref{chap:riepilogo_iii}.

\section{Il 6/4 cadenzale}
\label{sec:cadenziale}

\begin{definizione}[6/4 cadenzale]
    È l'accordo di \sigla{I}{6}{4} che precede immediatamente la dominante nella formula cadenzale. Nelle voci superiori ci sono la quarta e la sesta rispetto al basso, che sta sulla nota della dominante e vi resta fermo nell'accordo successivo.
\end{definizione}

\begin{regola}
    \begin{itemize}
        \item Le note in quarta e sesta devono essere \textbf{preparate} (tenute o raggiunte per grado congiunto).
        \item La quarta risolve scendendo alla terza della dominante.
        \item La sesta risolve scendendo alla quinta della dominante.
        \item Metricamente il 6/4 cade sul tempo \textbf{forte} e il V sul tempo successivo.
    \end{itemize}
\end{regola}

\begin{esempio}[Il 6/4 cadenzale in Do maggiore]
    Progressione IV - \sigla{I}{6}{4} - V - I.
    \musica{es_cap03_01}
\end{esempio}

\section{Il 6/4 di passaggio}
\label{sec:64_passaggio}

\begin{definizione}[6/4 di passaggio]
    Si forma quando il basso si muove per grado congiunto passando per la quinta dell'accordo in un tratto scalare.
\end{definizione}

\begin{regola}
    \begin{itemize}
        \item Il basso procede per grado congiunto sia in entrata sia in uscita.
        \item L'accordo cade su un tempo \textbf{debole}.
        \item Le voci superiori procedono per grado congiunto o restano ferme.
    \end{itemize}
\end{regola}

\begin{esempio}[6/4 di passaggio]
    Passaggio I - \sigla{V}{6}{4} - \sigla{I}{6}{} su basso scalare Do-Re-Mi: il contralto resta fermo, il soprano e il tenore procedono per grado congiunto.
    \musica{es_cap03_02}
\end{esempio}

\section{Il 6/4 di pedale o di ornamento}
\label{sec:64_pedale}

\begin{definizione}[6/4 di pedale]
    Si forma quando le voci superiori si muovono sopra un basso fermo (pedale).
\end{definizione}

\begin{regola}
    Il basso, fermo, appartiene all'accordo che precede e a quello che segue (di norma I - \sigla{IV}{6}{4} - I sulla tonica). Le voci superiori raggiungono e lasciano il 6/4 per grado congiunto: è un ornamento transitorio sopra una nota fissa.
\end{regola}
